% !TEX program = pdflatex
\documentclass[12pt,a4paper]{article}

% ─── Packages ───────────────────────────────────────────────────────────
\usepackage[utf8]{inputenc}
\usepackage[T1]{fontenc}
\usepackage{lmodern}
\usepackage[margin=2.5cm]{geometry}
\usepackage{amsmath,amssymb,amsfonts}
\usepackage{bm}
\usepackage{mathtools}
\usepackage{booktabs}
\usepackage{array}
\usepackage{graphicx}
\usepackage{float}
\usepackage{algorithm}
\usepackage{algpseudocode}
\usepackage{hyperref}
\usepackage{cleveref}
\usepackage{xcolor}
\usepackage{listings}
\usepackage{caption}
\usepackage{subcaption}
\usepackage{enumitem}
\usepackage{tikz}
\usetikzlibrary{arrows.meta,positioning,chains,calc}

\hypersetup{
    colorlinks=true,
    linkcolor=blue!70!black,
    citecolor=green!50!black,
    urlcolor=blue!80!black
}

% ─── Math operators ─────────────────────────────────────────────────────
\DeclareMathOperator{\RMSE}{RMSE}
\DeclareMathOperator{\svd}{svd}
\DeclareMathOperator{\diag}{diag}
\DeclareMathOperator{\rank}{rank}
\newcommand{\R}{\mathbb{R}}
\newcommand{\SE}{\mathrm{SE}}
\newcommand{\SO}{\mathrm{SO}}
\newcommand{\T}{\mathbf{T}}
\newcommand{\Rot}{\mathbf{R}}
\newcommand{\tvec}{\mathbf{t}}
\newcommand{\Rz}{\mathbf{R}_z}
\newcommand{\nh}{\hat{\mathbf{n}}}
\newcommand{\nraw}{\mathbf{n}}
\newcommand{\pvec}{\mathbf{p}}
\newcommand{\norm}[1]{\left\|#1\right\|}
\newcommand{\abs}[1]{\left|#1\right|}
\newcommand{\code}[1]{\texttt{#1}}

% ─── Document ───────────────────────────────────────────────────────────
\title{%
    \textbf{Kinematic Calibration of a Pan-Tilt Mounted\\
    Planar LiDAR System}\\[0.5em]
    \large Mathematical Formulation, Implementation, and Experimental Results
}
\author{%
    Mert \\
    Bitirme Projesi (Graduation Project)
}
\date{March 2026}

\begin{document}
\maketitle

\begin{abstract}
This report presents a comprehensive kinematic calibration methodology for a
pan-tilt unit (PTU) carrying a Hokuyo URG-04LX planar LiDAR sensor. The system
sweeps the 2D laser across multiple pan-tilt configurations while facing a
flat wall, and estimates 14 kinematic correction parameters by minimizing the
global plane-consistency cost of the projected 3D points. We detail the
complete mathematical formulation of every stage of the pipeline:
grid-scan data acquisition with RANSAC line fitting, per-capture line quality
filtering, forward kinematics with correction parameterization,
least-squares optimization with Tikhonov regularization, observability
analysis via numerical SVD, and cross-validation. Experimental results on
31 scan configurations demonstrate a 12.3\% reduction in global plane RMSE
after calibration, with the per-capture line quality filter alone accounting
for a 40\% noise reduction. We identify fundamental observability limitations
of single-wall calibration and discuss strategies for achieving higher accuracy.
\end{abstract}

\tableofcontents
\newpage

% ═══════════════════════════════════════════════════════════════════════
\section{Introduction}
\label{sec:intro}
% ═══════════════════════════════════════════════════════════════════════

A planar LiDAR sensor mounted on a pan-tilt mechanism can generate
three-dimensional point clouds by sweeping the 2D scan plane through a
range of orientations. The accuracy of the resulting 3D reconstruction
depends critically on the precise knowledge of the kinematic chain
relating the laser frame to a fixed base frame. Nominal values from
hand measurements and mechanical drawings are typically insufficient
for applications requiring sub-centimetre accuracy.

\textbf{Kinematic calibration} estimates the true geometric parameters of the
mechanism---joint origins, orientations, and encoder offsets---from observed
data. The fundamental idea is to view a geometrically known target (a flat
wall) from multiple configurations and solve for the parameters that best
explain the observations.

\subsection{System Overview}

Our system consists of:
\begin{itemize}[nosep]
    \item \textbf{LiDAR:} Hokuyo URG-04LX, planar scanner,
          240\textdegree{} FOV, $\sim$60\,mm accuracy specification.
    \item \textbf{Pan-Tilt Unit:} Two-axis servo mechanism controlled via
          Arduino, reporting joint states via ROS \code{/joint\_states}.
    \item \textbf{Software:} ROS~1 Noetic, Python 3 calibration pipeline.
\end{itemize}

\subsection{Report Outline}

The remainder of this report is organized as follows:
\Cref{sec:kinematics} describes the nominal kinematic chain from the URDF
and the correction parameterization.
\Cref{sec:acquisition} details the data acquisition procedure including
grid generation, FOV extraction, and RANSAC line fitting.
\Cref{sec:filter} presents the per-capture line quality filter.
\Cref{sec:solver} develops the optimization formulation, cost function,
and solver configuration.
\Cref{sec:observability} covers the observability analysis.
\Cref{sec:validation} describes the validation methodology.
\Cref{sec:results} presents the experimental results.
\Cref{sec:future} discusses limitations and future work needed for
improved calibration.


% ═══════════════════════════════════════════════════════════════════════
\section{Kinematic Chain and Correction Parameterization}
\label{sec:kinematics}
% ═══════════════════════════════════════════════════════════════════════

\subsection{URDF Kinematic Chain}

The mechanism is described by the following URDF link--joint chain:
\[
    \texttt{base\_link}
    \xrightarrow{\text{joint2 (pan)}}
    \texttt{pan\_link}
    \xrightarrow{\text{joint1 (tilt)}}
    \texttt{tilt\_link}
    \xrightarrow{\text{hokuyo\_base\_joint}}
    \texttt{hokuyo\_base}
    \xrightarrow{\text{hokuyo\_joint}}
    \texttt{laser}
\]

Each joint has a fixed origin transform $\T_{\text{origin}} \in \SE(3)$
specified by a translation vector $\mathbf{t} = (x, y, z)$ and an
orientation given by roll-pitch-yaw angles $(\phi, \theta, \psi)$.

\subsubsection{Homogeneous Transform from URDF Parameters}

Given $\mathbf{t} = (x,y,z)$ and $\text{rpy} = (\phi, \theta, \psi)$,
the URDF convention constructs the $4 \times 4$ homogeneous transformation
matrix as:
\begin{equation}
    \T(\mathbf{t}, \text{rpy}) =
    \begin{bmatrix}
        \Rot_z(\psi)\,\Rot_y(\theta)\,\Rot_x(\phi) & \mathbf{t} \\
        \mathbf{0}^\top & 1
    \end{bmatrix}
    \label{eq:urdf_transform}
\end{equation}
where $\Rot_x(\phi)$, $\Rot_y(\theta)$, $\Rot_z(\psi)$ are the elementary
rotation matrices about the $x$, $y$, $z$ axes respectively.

\textbf{Critical convention:} The product $\Rot_z(\psi)\,\Rot_y(\theta)\,\Rot_x(\phi)$
corresponds to \emph{extrinsic} $X$-$Y$-$Z$ rotations (equivalently,
\emph{intrinsic} $Z$-$Y'$-$X''$). In the \code{scipy} library, this is
specified as \code{Rotation.from\_euler('XYZ', rpy)}, where the uppercase
letters denote the extrinsic (fixed-axis) convention.

\begin{center}
\fbox{%
\parbox{0.85\textwidth}{%
\textbf{Bug found during this project:}
The original code used \code{'xyz'} (lowercase, intrinsic),
which applies rotations in the opposite order. Correcting to
\code{'XYZ'} (extrinsic) reduced the uncalibrated plane RMSE
from 56\,mm to 27\,mm---a single-character fix with enormous impact.}}
\end{center}

\subsubsection{Nominal Transform Values}

The nominal (hand-measured) URDF parameters are:

\begin{table}[H]
\centering
\caption{Nominal URDF joint parameters.}
\label{tab:nominal_urdf}
\begin{tabular}{@{}llccc@{}}
\toprule
Joint & Frames & xyz (m) & rpy (rad) & Axis \\
\midrule
\code{joint2} (pan) & base $\to$ pan
    & $(0,\ {-0.0285},\ 0.092)$ & $(0,\, 0,\, 0)$ & $(0,\,0,\,{-1})$ \\
\code{joint1} (tilt) & pan $\to$ tilt
    & $(0,\ {-0.009},\ 0.022)$ & $(0,\, 1.57,\, 0)$ & $(0,\,0,\,1)$ \\
\code{hokuyo\_base} & tilt $\to$ hokuyo\_base
    & $(-0.045,\, 0,\, 0)$ & $(0,\, 0,\, 0)$ & --- \\
\code{hokuyo\_joint} & hokuyo\_base $\to$ laser
    & $(-0.061,\, 0.014,\, 0.040)$ & $(1.57,\, 3.14,\, {-1.57})$ & --- \\
\bottomrule
\end{tabular}
\end{table}

In our implementation, the two fixed joints (\code{hokuyo\_base\_joint}
and \code{hokuyo\_joint}) are composed into a single transform:
\begin{equation}
    \T_{\text{laser}}^{\text{nom}} = \T_{\text{hbase}} \cdot \T_{\text{hjoint}}
    \label{eq:T_laser_nominal}
\end{equation}

The nominal static transforms used in forward kinematics are therefore:
\begin{align}
    \T_{\text{pan}}^{\text{nom}} &= \T\bigl((0,\, {-0.0285},\, 0.092),\; (0,\,0,\,0)\bigr)
    \label{eq:T_pan_nom} \\
    \T_{\text{tilt}}^{\text{nom}} &= \T\bigl((0,\, {-0.009},\, 0.022),\; (0,\, 1.57,\, 0)\bigr)
    \label{eq:T_tilt_nom} \\
    \T_{\text{laser}}^{\text{nom}} &= \T\bigl(({-0.045},\, 0,\, 0),\; (0,\,0,\,0)\bigr)
    \cdot \T\bigl(({-0.061},\, 0.014,\, 0.040),\; (1.57,\, 3.14,\, {-1.57})\bigr)
    \label{eq:T_laser_nom}
\end{align}

\subsection{Correction Parameterization}
\label{sec:correction_param}

We parameterize kinematic errors as small corrections applied to the
nominal transforms. Consider a nominal transform $\T^{\text{nom}}$ and
a small perturbation consisting of a translation $\delta\mathbf{t} \in \R^3$
and a rotation $\delta\mathbf{r} = (\delta\phi, \delta\theta, \delta\psi) \in \R^3$:
\begin{equation}
    \T^{\text{corr}} = \T^{\text{nom}} \cdot \Delta\T, \qquad
    \Delta\T =
    \begin{bmatrix}
        \Rot(\delta\mathbf{r}) & \delta\mathbf{t} \\
        \mathbf{0}^\top & 1
    \end{bmatrix}
    \label{eq:correction}
\end{equation}
where $\Rot(\delta\mathbf{r}) = \Rot_z(\delta\psi)\,\Rot_y(\delta\theta)\,\Rot_x(\delta\phi)$
follows the same extrinsic $XYZ$ convention. This right-multiplication
means the correction is expressed in the \emph{child} frame of the joint,
which is the natural choice for body-frame perturbations.

\subsection{The 14-Parameter Model}
\label{sec:14param}

For calibration with a single flat wall, the tilt joint origin corrections
are \emph{degenerate} with the laser mount corrections (see \Cref{sec:observability}).
We therefore fix the tilt joint geometry at its nominal values and define
a 14-parameter correction vector:
\begin{equation}
    \bm{\delta} =
    \bigl[\underbrace{\delta\mathbf{t}_{\text{pan}},\; \delta\mathbf{r}_{\text{pan}}}_{\text{6: pan origin}},\;
    \underbrace{\varepsilon_{\text{pan}},\; \varepsilon_{\text{tilt}}}_{\text{2: encoder zeros}},\;
    \underbrace{\delta\mathbf{t}_{\text{laser}},\; \delta\mathbf{r}_{\text{laser}}}_{\text{6: laser mount}}\bigr]
    \in \R^{14}
    \label{eq:delta}
\end{equation}

The parameters are:
\begin{enumerate}[nosep]
    \item $\delta\mathbf{t}_{\text{pan}} = (\delta t_x, \delta t_y, \delta t_z)$:
          Translation correction to the pan joint origin.
    \item $\delta\mathbf{r}_{\text{pan}} = (\delta\phi, \delta\theta, \delta\psi)$:
          Rotation correction to the pan joint origin.
    \item $\varepsilon_{\text{pan}}$: Pan encoder zero offset (rad). The true pan angle
          is $\theta_{\text{pan}} + \varepsilon_{\text{pan}}$.
    \item $\varepsilon_{\text{tilt}}$: Tilt encoder zero offset (rad). The true tilt angle
          is $\theta_{\text{tilt}} + \varepsilon_{\text{tilt}}$.
    \item $\delta\mathbf{t}_{\text{laser}}$: Translation correction to the composite
          laser transform.
    \item $\delta\mathbf{r}_{\text{laser}}$: Rotation correction to the composite
          laser transform.
\end{enumerate}

\subsection{Forward Kinematics with Corrections}
\label{sec:fk}

Given joint angles $(\theta_p, \theta_t)$ and correction vector $\bm{\delta}$,
the forward kinematics computes the laser-to-base transform:

\begin{equation}
\boxed{
    \T_{\text{laser}}^{\text{base}}(\theta_p, \theta_t;\, \bm{\delta}) =
    \T_{\text{pan}}^{\text{corr}} \cdot
    \Rz\!\bigl({-(\theta_p + \varepsilon_{\text{pan}})}\bigr) \cdot
    \T_{\text{tilt}}^{\text{nom}} \cdot
    \Rz\!\bigl(\theta_t + \varepsilon_{\text{tilt}}\bigr) \cdot
    \T_{\text{laser}}^{\text{corr}}
}
\label{eq:fk}
\end{equation}

where:
\begin{align}
    \T_{\text{pan}}^{\text{corr}} &= \T_{\text{pan}}^{\text{nom}} \cdot \Delta\T(\delta\mathbf{t}_{\text{pan}}, \delta\mathbf{r}_{\text{pan}})
    \label{eq:T_pan_corr} \\
    \T_{\text{laser}}^{\text{corr}} &= \T_{\text{laser}}^{\text{nom}} \cdot \Delta\T(\delta\mathbf{t}_{\text{laser}}, \delta\mathbf{r}_{\text{laser}})
    \label{eq:T_laser_corr}
\end{align}
and $\Rz(\alpha)$ is the $4 \times 4$ homogeneous rotation about $z$:
\begin{equation}
    \Rz(\alpha) =
    \begin{bmatrix}
        \cos\alpha & -\sin\alpha & 0 & 0 \\
        \sin\alpha &  \cos\alpha & 0 & 0 \\
        0          &  0          & 1 & 0 \\
        0          &  0          & 0 & 1
    \end{bmatrix}
    \label{eq:Rz}
\end{equation}

\textbf{Joint axis conventions:}
\begin{itemize}[nosep]
    \item The pan joint (joint2) has axis $(0,\, 0,\, {-1})$, so the rotation
          is $\Rz(-\theta)$ rather than $\Rz(\theta)$.
    \item The tilt joint (joint1) has axis $(0,\, 0,\, 1)$. Its origin includes
          an RPY of $(0,\, 1.57,\, 0)$ which rotates the local $z$-axis to
          effectively create tilt motion.
\end{itemize}

\subsection{Point Projection}

A 2D laser point $\pvec_L = (x_L,\, y_L)$ in the laser frame
is projected to the base frame via:
\begin{equation}
    \pvec_{\text{base}} = \T_{\text{laser}}^{\text{base}} \cdot
    \begin{bmatrix} x_L \\ y_L \\ 0 \\ 1 \end{bmatrix}
    \label{eq:point_projection}
\end{equation}
The third coordinate is zero because the planar LiDAR scans in the $xy$-plane
of its own frame.


% ═══════════════════════════════════════════════════════════════════════
\section{Data Acquisition}
\label{sec:acquisition}
% ═══════════════════════════════════════════════════════════════════════

\subsection{Grid-Scan Strategy}

We sample the joint space on a regular grid:
\begin{equation}
    \mathcal{G} = \bigl\{(\theta_p^{(i)}, \theta_t^{(j)}) :
    \theta_p^{(i)} \in \text{linspace}(-\Theta_p, +\Theta_p, N_p),\;
    \theta_t^{(j)} \in \text{linspace}(-\Theta_t, +\Theta_t, N_t)\bigr\}
    \label{eq:grid}
\end{equation}
where $\Theta_p$ and $\Theta_t$ are the pan and tilt limits (typically 25\textdegree{}
and 15\textdegree{}), and $N_p = N_t = 5$ gives a $5 \times 5 = 25$
configuration grid (expanded to $5 \times 7 = 35$ in practice with
intermediate tilt values).

At each grid pose $(\theta_p^{(k)}, \theta_t^{(k)})$:
\begin{enumerate}[nosep]
    \item The PTU is commanded to the target angles via
          \code{/joint\_command} at 50\,Hz.
    \item Convergence is verified by checking
          $|\theta_{\text{current}} - \theta_{\text{target}}| < \epsilon$
          where $\epsilon = 1$\textdegree.
    \item After convergence, the system waits a settle time $t_s = 1.5$\,s.
    \item A single \code{/scan} message is captured.
    \item The actual encoder-reported joint angles are recorded from
          \code{/joint\_states}.
\end{enumerate}

\subsection{FOV Extraction}
\label{sec:fov}

The Hokuyo URG-04LX has a 240\textdegree{} field of view, but the
calibration wall subtends only a small angular region near the forward
direction. We extract a region of interest (ROI) by keeping only
scan points whose bearing angle satisfies:
\begin{equation}
    \tilde{P}_k = \bigl\{ \pvec \in P_k : |\alpha(\pvec)| < \alpha_{\max} \bigr\}
    \label{eq:fov_extraction}
\end{equation}
where $\alpha(\pvec)$ is the bearing of point $\pvec$ in the laser frame
and $\alpha_{\max}$ is the half-FOV parameter (typically 10--15\textdegree).
Only points within the valid range interval $[r_{\min},\, r_{\max}]$ are retained.

Each valid range measurement $r_i$ at angle $\alpha_i = \alpha_{\min} + i \cdot \Delta\alpha$
is converted to Cartesian coordinates:
\begin{equation}
    \pvec_i = \begin{pmatrix} r_i \cos\alpha_i \\ r_i \sin\alpha_i \end{pmatrix}
    \label{eq:polar_to_cart}
\end{equation}

\subsection{RANSAC Line Fitting}
\label{sec:ransac}

Since the laser is viewing a flat wall, the extracted 2D points should
lie on a straight line. We use RANSAC (Random Sample Consensus) to
robustly fit this line and reject outliers.

\begin{algorithm}[H]
\caption{RANSAC 2D Line Fit}
\label{alg:ransac}
\begin{algorithmic}[1]
\Require Point set $\tilde{P} = \{\pvec_1, \ldots, \pvec_N\} \subset \R^2$,
         threshold $\tau$, iterations $K$
\Ensure Line direction $\hat{\mathbf{v}}$, mean $\bar{\pvec}$, inlier set $\mathcal{I}$
\State $n_{\text{best}} \gets 0$
\For{$k = 1, \ldots, K$} \Comment{$K = 600$}
    \State Sample two distinct points $\pvec_a, \pvec_b$ uniformly at random
    \State $\mathbf{v} \gets (\pvec_b - \pvec_a) / \norm{\pvec_b - \pvec_a}$
    \State $\mathbf{n} \gets (-v_y,\; v_x)$ \Comment{Line normal (perpendicular)}
    \State $d_i \gets |(\pvec_i - \pvec_a)^\top \mathbf{n}|$ for each $\pvec_i$
    \State $n_{\text{inliers}} \gets |\{i : d_i \le \tau\}|$
    \Comment{$\tau = 15$\,mm}
    \If{$n_{\text{inliers}} > n_{\text{best}}$}
        \State Update best model
    \EndIf
    \If{$n_{\text{inliers}} > 0.7 \cdot N$}
        \State \textbf{break} \Comment{Early termination}
    \EndIf
\EndFor
\State Collect inliers $\mathcal{I}$ for the best model
\State Refine: compute $\bar{\pvec} = \text{mean}(\mathcal{I})$, centre inliers,
       compute SVD to get refined direction $\hat{\mathbf{v}}$ as the first
       right singular vector of the centred inlier matrix
\end{algorithmic}
\end{algorithm}

The SVD refinement step computes the best-fit line through the inliers
in a least-squares sense. Given the $M \times 2$ centred inlier matrix
$\mathbf{C} = \mathcal{I} - \mathbf{1}\bar{\pvec}^\top$, we compute
$\mathbf{C} = \mathbf{U}\bm{\Sigma}\mathbf{V}^\top$ and take
$\hat{\mathbf{v}} = \mathbf{V}_{:,1}$ (the direction of maximum variance).

Each capture stores: inlier points $\mathcal{I}_k$, line direction
$\hat{\mathbf{v}}_k$, line mean $\bar{\pvec}_k$, pan angle $\theta_p^{(k)}$,
tilt angle $\theta_t^{(k)}$, and timestamp.

\subsection{Data Storage}

All captures are saved to a compressed NumPy archive (\code{.npz}) containing:
\begin{itemize}[nosep]
    \item \code{n\_captures}: number of successful captures
    \item \code{fov\_half\_deg}: the FOV parameter used
    \item For each capture $k$: \code{cap\_\textit{k}\_pan\_rad},
          \code{cap\_\textit{k}\_tilt\_rad}, \code{cap\_\textit{k}\_line\_dir},
          \code{cap\_\textit{k}\_line\_mean}, \code{cap\_\textit{k}\_inlier\_points},
          \code{cap\_\textit{k}\_all\_points}, \code{cap\_\textit{k}\_timestamp}
\end{itemize}


% ═══════════════════════════════════════════════════════════════════════
\section{Per-Capture Line Quality Filter}
\label{sec:filter}
% ═══════════════════════════════════════════════════════════════════════

Before the calibration solver processes the data, we apply a quality
filter based on the following observation:

\begin{center}
\fbox{\parbox{0.85\textwidth}{%
\textbf{Principle:} If a single scan's inlier points do not conform to a
straight line with sufficient precision, those points cannot reliably
contribute to a global plane fit. Since the Hokuyo's scan plane intersects
a flat wall in a line, the per-capture line RMSE is a lower bound on
the achievable 3D reconstruction accuracy.}}
\end{center}

\subsection{Line RMSE Computation}

For capture $k$ with inlier points $\mathcal{I}_k$, fitted line direction
$\hat{\mathbf{v}}_k$, and line mean $\bar{\pvec}_k$, the
perpendicular distance of each point to the line is:
\begin{equation}
    d_i = (\pvec_i - \bar{\pvec}_k)^\top \mathbf{n}_\perp, \qquad
    \mathbf{n}_\perp = \begin{pmatrix} -\hat{v}_y \\ \hat{v}_x \end{pmatrix}
    \label{eq:line_dist}
\end{equation}
where $\mathbf{n}_\perp$ is the unit normal perpendicular to the line direction.

The line RMSE for capture $k$ is:
\begin{equation}
    \text{RMSE}_{\text{line}}^{(k)} = \sqrt{\frac{1}{|\mathcal{I}_k|}
    \sum_{i \in \mathcal{I}_k} d_i^2}
    \label{eq:line_rmse}
\end{equation}

\subsection{Two-Stage Filtering}

The filter operates in two stages:

\paragraph{Stage 1: Capture Rejection.}
A capture is rejected entirely if its line RMSE exceeds a threshold $\tau_{\text{line}}$:
\begin{equation}
    \text{RMSE}_{\text{line}}^{(k)} > \tau_{\text{line}} \implies \text{reject capture } k
    \label{eq:capture_reject}
\end{equation}
The default threshold is $\tau_{\text{line}} = 3.0$\,mm.

\paragraph{Stage 2: Point Clipping.}
Within accepted captures, individual points with excessive perpendicular
error are removed:
\begin{equation}
    |d_i| \cdot 1000 > \tau_{\text{pt}} \implies \text{remove point } i
    \label{eq:point_clip}
\end{equation}
where $\tau_{\text{pt}} = 2 \cdot \tau_{\text{line}}$ by default (6.0\,mm).

\subsection{Empirical Impact}

In our dataset of 31 captures, applying the 3.0\,mm line quality filter:
\begin{itemize}[nosep]
    \item Rejected 12 captures (mostly at extreme pan angles $\pm 25$\textdegree{}
          and $+15$\textdegree), retaining 19.
    \item Clipped 12 individual outlier points from the remaining captures.
    \item Retained 1012 of 1612 original points.
    \item \textbf{Reduced uncalibrated plane RMSE from 26.9\,mm to 16.2\,mm}
          ---a 40\% improvement from filtering alone.
\end{itemize}

The captures at extreme pan angles exhibit higher noise because the
wall is seen at a shallow angle, increasing the range uncertainty's
effect on the reconstructed point position.


% ═══════════════════════════════════════════════════════════════════════
\section{Calibration Solver}
\label{sec:solver}
% ═══════════════════════════════════════════════════════════════════════

\subsection{Point Subsampling}

To control computational cost, inlier points within each capture are
deterministically subsampled:
\begin{equation}
    \text{step} = \max\Bigl(1,\; \Bigl\lfloor \frac{|\mathcal{I}_k|}{M_{\max}} \Bigr\rfloor\Bigr),
    \qquad
    \mathcal{I}_k^{\text{sub}} = \{\pvec_{0},\, \pvec_{\text{step}},\, \pvec_{2\cdot\text{step}},\, \ldots\}
    \label{eq:subsample}
\end{equation}
where $M_{\max} = 30$ points per capture. This ensures uniform spatial
coverage within each scan while keeping the total point count manageable.

\subsection{Initial Plane Estimation}
\label{sec:init_plane}

Before optimization, we estimate the wall plane from uncalibrated
forward kinematics ($\bm{\delta} = \mathbf{0}$). All subsampled points are
projected to the base frame and stacked into a matrix $\mathbf{P} \in \R^{M \times 3}$.
The plane is fitted via SVD:
\begin{enumerate}[nosep]
    \item Compute the centroid: $\bar{\pvec} = \frac{1}{M}\sum_{i=1}^{M} \pvec_i$
    \item Centre the points: $\tilde{\mathbf{P}} = \mathbf{P} - \mathbf{1}\bar{\pvec}^\top$
    \item Compute the SVD: $\tilde{\mathbf{P}} = \mathbf{U}\bm{\Sigma}\mathbf{V}^\top$
    \item The plane normal is the third right singular vector:
          $\nh = \mathbf{V}_{:,3}$ (direction of minimum variance)
    \item The plane offset is: $d = \nh^\top \bar{\pvec}$
\end{enumerate}

This yields the initial plane $(\nh_0, d_0)$ used to seed the optimizer.
Its sign is chosen so that $d_0 > 0$ (wall is in front of the sensor).

\subsection{Optimization Variables}

The full optimization variable is:
\begin{equation}
    \mathbf{x} = \bigl[\underbrace{\bm{\delta}}_{\text{14 kinematic}},\;
    \underbrace{\nraw,\; d}_{\text{4 plane}}\bigr] \in \R^{18}
    \label{eq:opt_vars}
\end{equation}
The plane parameters $(\nraw, d)$ are co-optimized because the ``true'' wall
plane is unknown and must be estimated simultaneously with the kinematic
corrections. This avoids the chicken-and-egg problem of needing the plane
to evaluate the cost, but needing good kinematics to estimate the plane.

\subsection{Cost Function}
\label{sec:cost}

The cost function is a sum of residuals
$\mathbf{r}(\mathbf{x}) = [r_1, \ldots, r_N]$ minimized in the
least-squares sense $\min_{\mathbf{x}} \frac{1}{2}\norm{\mathbf{r}(\mathbf{x})}^2$.
The residual vector consists of three components:

\subsubsection{(A) Global Plane Consistency}

For each capture $k$ and each subsampled point $\pvec_j \in \mathcal{I}_k^{\text{sub}}$:
\begin{equation}
    r_{k,j}^{\text{plane}} = \alpha \cdot \bigl(\nh^\top \pvec_{j}^{\text{base}} - d\bigr)
    \label{eq:plane_residual}
\end{equation}
where $\pvec_j^{\text{base}} = \T_{\text{laser}}^{\text{base}}(\theta_p^{(k)}, \theta_t^{(k)}; \bm{\delta}) \cdot (\pvec_j^L, 0, 1)^\top$
is the point projected to the base frame, $\nh = \nraw / \norm{\nraw}$
is the unit normal, and $\alpha = 5.0$ is the plane weight.

The quantity $\nh^\top \pvec - d$ is the \emph{signed distance} from
point $\pvec$ to the plane $\{\mathbf{x} : \nh^\top \mathbf{x} = d\}$.
For a perfect calibration and a truly flat wall, all these residuals
would be zero.

\subsubsection{(B) Normal Unit-Length Constraint}

Since $\nraw$ is optimized as a free 3-vector, we add a soft constraint
to keep it unit-length:
\begin{equation}
    r^{\text{norm}} = \omega_n \cdot \bigl(\norm{\nraw} - 1\bigr),
    \qquad \omega_n = 10
    \label{eq:norm_constraint}
\end{equation}
This penalty drives $\norm{\nraw} \to 1$ without requiring an explicit
constraint, which would complicate the least-squares formulation.

\subsubsection{(C) Tikhonov Regularization}

To prevent overfitting to noise (particularly for poorly observed
parameters), we add an $\ell_2$ penalty on the correction vector:
\begin{equation}
    r_i^{\text{reg}} = \sqrt{\lambda} \cdot \delta_i, \qquad i = 1, \ldots, 14
    \label{eq:regularization}
\end{equation}
where $\lambda$ is the regularization weight. This is equivalent to the
Tikhonov regularization term $\lambda \norm{\bm{\delta}}^2$ in the
squared cost, which penalizes large deviations from the nominal geometry
and implicitly encodes the prior that the hand-measured URDF is
approximately correct.

\subsubsection{Complete Residual Vector}

The full residual vector has $N_{\text{pts}} + 1 + 14$ elements:
\begin{equation}
    \mathbf{r}(\mathbf{x}) =
    \begin{bmatrix}
        \alpha \cdot (\nh^\top \pvec_1^{\text{base}} - d) \\
        \vdots \\
        \alpha \cdot (\nh^\top \pvec_{N_{\text{pts}}}^{\text{base}} - d) \\
        10 \cdot (\norm{\nraw} - 1) \\
        \sqrt{\lambda} \cdot \delta_1 \\
        \vdots \\
        \sqrt{\lambda} \cdot \delta_{14}
    \end{bmatrix}
    \label{eq:full_residuals}
\end{equation}

\subsection{Solver Configuration}

The optimization problem is solved using \code{scipy.optimize.least\_squares}
with the Trust Region Reflective (\code{trf}) method, which supports
box constraints. The bounds are:
\begin{equation}
    \begin{aligned}
        |\delta t_i| &\le B_t  &&\text{(translation bound, default 5\,mm)} \\
        |\delta r_i| &\le B_r  &&\text{(rotation bound, default 0.05\,rad $\approx$ 2.87\textdegree)} \\
        |\varepsilon_{\text{enc}}| &\le B_e  &&\text{(encoder bound, default 10\textdegree)} \\
        |n_i| &\le 1,\quad |d| \le 20  &&\text{(plane parameters)}
    \end{aligned}
    \label{eq:bounds}
\end{equation}

The solver uses very tight convergence criteria (\code{ftol} = \code{xtol}
= \code{gtol} = $10^{-12}$) to ensure the solution is driven by
bound constraints or the cost landscape rather than premature termination.

\subsection{Output}

The solver returns:
\begin{itemize}[nosep]
    \item The optimal correction vector $\bm{\delta}^* \in \R^{14}$.
    \item Encoder zero offsets $(\varepsilon_{\text{pan}}^*, \varepsilon_{\text{tilt}}^*)$.
    \item Corrected URDF joint origins computed by applying $\bm{\delta}^*$
          to the nominal transforms via \cref{eq:correction}.
\end{itemize}


% ═══════════════════════════════════════════════════════════════════════
\section{Observability Analysis}
\label{sec:observability}
% ═══════════════════════════════════════════════════════════════════════

Not all 14 kinematic parameters may be identifiable from the available
data. The \emph{observability analysis} determines which parameters (or
linear combinations thereof) can actually be estimated.

\subsection{Numerical Jacobian}

We compute the Jacobian of the residual vector with respect to the
14 kinematic parameters using finite differences:
\begin{equation}
    J_{i,j} = \frac{r_i(\mathbf{x} + \epsilon \mathbf{e}_j) - r_i(\mathbf{x})}
    {\epsilon}, \qquad j = 1, \ldots, 14
    \label{eq:jacobian}
\end{equation}
where $\epsilon = 10^{-7}$ and $\mathbf{e}_j$ is the $j$-th standard basis
vector (applied only to the first 14 components of $\mathbf{x}$).

\subsection{SVD-Based Rank Analysis}

The Jacobian $\mathbf{J} \in \R^{N_r \times 14}$ is decomposed via SVD:
\begin{equation}
    \mathbf{J} = \mathbf{U} \bm{\Sigma} \mathbf{V}^\top
    \label{eq:svd}
\end{equation}
where $\bm{\Sigma} = \diag(\sigma_1, \ldots, \sigma_{14})$ with
$\sigma_1 \ge \sigma_2 \ge \cdots \ge \sigma_{14} \ge 0$.

The \emph{observable rank} is the number of singular values exceeding a
threshold $\sigma_{\min}$:
\begin{equation}
    r = \bigl|\{k : \sigma_k > 1.5\sqrt{\lambda}\}\bigr|
    \label{eq:obs_rank}
\end{equation}
The threshold $1.5\sqrt{\lambda}$ is chosen to be slightly above the
Tikhonov regularization floor, so that parameters whose identifiability
comes only from the regularization term are excluded.

\subsection{Per-Parameter Observability Score}

Each parameter is assigned an observability score based on its
contribution to the well-observed singular directions:
\begin{equation}
    s_j = \frac{1}{s_{\max}} \sum_{k:\, \sigma_k > \sigma_{\min}}
    \sigma_k^2 \cdot V_{k,j}^2
    \label{eq:obs_score}
\end{equation}
where $s_{\max} = \max_j s_j$ normalizes the scores to $[0, 1]$.
Parameters with $s_j > 0.05$ are marked as \emph{well-observed} (V),
while others are marked as \emph{poorly observed} (X).

\subsection{Interpretation for Single-Wall Calibration}

With our $5 \times 7$ grid of captures facing a single flat wall,
the analysis consistently identifies:
\begin{itemize}[nosep]
    \item \textbf{Observable rank: 7 out of 14.}
    \item Well-observed parameters: a subset of pan and laser translations.
    \item Poorly observed (null space): all pan rotations, encoder zeros,
          laser $z$-translation, and all laser rotations.
\end{itemize}

The 7 unobservable directions correspond to parameter combinations that
produce identical projected point patterns on a single flat surface.
This is a fundamental geometric limitation, not a software issue.


% ═══════════════════════════════════════════════════════════════════════
\section{Validation}
\label{sec:validation}
% ═══════════════════════════════════════════════════════════════════════

Validation measures how well the calibrated kinematic chain explains
held-out (or same-set) observations, compared to the uncalibrated chain.

\subsection{Validation Procedure}

Given a validation dataset of $K$ captures and the optimal correction
$\bm{\delta}^*$:

\begin{enumerate}[nosep]
    \item \textbf{Apply line quality filter} (same threshold as solver).
    \item For each capture $k$, transform inlier points to base frame:
          $\pvec_j^{\text{base}} = \T_{\text{laser}}^{\text{base}}(\theta_p^{(k)}, \theta_t^{(k)}; \bm{\delta}^*) \cdot (\pvec_j^L, 0, 1)^\top$
    \item Fit a global plane $(\nh, d)$ via SVD (\Cref{sec:init_plane}).
    \item Compute per-capture plane RMSE:
          \begin{equation}
              \text{RMSE}_k = \sqrt{\frac{1}{|\mathcal{I}_k|} \sum_{j \in \mathcal{I}_k}
              (\nh^\top \pvec_j^{\text{base}} - d)^2} \times 1000
              \qquad \text{[mm]}
              \label{eq:val_rmse}
          \end{equation}
    \item Repeat steps 2--4 with $\bm{\delta} = \mathbf{0}$
          (uncalibrated) for comparison.
\end{enumerate}

\subsection{Reported Metrics}

\begin{itemize}[nosep]
    \item \textbf{Global RMSE:}
          $\text{RMSE}_{\text{global}} = \sqrt{\frac{1}{N}\sum_{i=1}^{N} e_i^2}$
          where $e_i$ are all point-to-plane signed distances across all captures.
    \item \textbf{Per-capture RMSE mean / std / max.}
    \item \textbf{Max point error:} $\max_i |e_i|$.
    \item \textbf{Improvement:} $(1 - \text{RMSE}_{\text{cal}} / \text{RMSE}_{\text{uncal}}) \times 100\%$.
    \item \textbf{PASS/FAIL verdict:} Based on whether the global RMSE falls
          below a user-defined threshold (default: 10\,mm).
\end{itemize}

\textbf{Important:} The line quality filter must be applied consistently
in both the solver and validation. Applying it only in the solver but
evaluating on unfiltered data leads to misleading results (calibrated
appearing worse than uncalibrated) because defect corrections amplify
noise in poor-quality captures.


% ═══════════════════════════════════════════════════════════════════════
\section{Experimental Results}
\label{sec:results}
% ═══════════════════════════════════════════════════════════════════════

\subsection{Dataset}

\begin{table}[H]
\centering
\caption{Calibration dataset characteristics.}
\label{tab:dataset}
\begin{tabular}{@{}ll@{}}
\toprule
Property & Value \\
\midrule
Total captures & 31 \\
Pan range & $-25$\textdegree{} to $+25$\textdegree{} \\
Tilt range & $-15$\textdegree{} to $+15$\textdegree{} \\
FOV half-angle & 10\textdegree{} \\
Total inlier points & 1612 \\
After line filter (3\,mm) & 19 captures, 1012 points \\
Subsampled for solver & 570 points \\
\bottomrule
\end{tabular}
\end{table}

\subsection{Euler Convention Bug}

The most impactful finding was the Euler convention mismatch.
\Cref{tab:euler_bug} shows the effect of correcting from intrinsic
\code{'xyz'} to extrinsic \code{'XYZ'}:

\begin{table}[H]
\centering
\caption{Impact of correcting the Euler convention.}
\label{tab:euler_bug}
\begin{tabular}{@{}lcc@{}}
\toprule
 & \code{'xyz'} (intrinsic, wrong) & \code{'XYZ'} (extrinsic, correct) \\
\midrule
Uncalibrated plane RMSE & 56\,mm & 27\,mm \\
\bottomrule
\end{tabular}
\end{table}

\subsection{Line Quality Filter Impact}

\begin{table}[H]
\centering
\caption{Effect of line quality filter on uncalibrated RMSE.}
\label{tab:filter_impact}
\begin{tabular}{@{}lcc@{}}
\toprule
 & Unfiltered (31 cap) & Filtered (19 cap) \\
\midrule
Uncalibrated plane RMSE & 26.9\,mm & 16.2\,mm \\
Reduction & --- & 40\% \\
\bottomrule
\end{tabular}
\end{table}

\subsection{Solver and Validation Results}

\begin{table}[H]
\centering
\caption{Calibration results with different solver settings
(all with 3\,mm line filter).}
\label{tab:solver_results}
\begin{tabular}{@{}lccccc@{}}
\toprule
Settings & $\lambda$ & $B_t$ / $B_r$ / $B_e$ & Obs.\ Rank
    & Cost $\downarrow$ & Global RMSE $\downarrow$ \\
\midrule
Tight bounds
    & 0.1 & 5\,mm / 2.87\textdegree{} / 3\textdegree{}
    & 7/14 & 20.1\% & 10.1\% \\
Wide encoder
    & 0.1 & 5\,mm / 2.87\textdegree{} / 10\textdegree{}
    & 7/14 & 22.3\% & 12.3\% \\
Strong reg.
    & 1.0 & 5\,mm / 2.87\textdegree{} / 10\textdegree{}
    & 6/14 & 21.4\% & 12.3\% \\
\bottomrule
\end{tabular}
\end{table}

The best configuration achieves a \textbf{12.3\% reduction} in global
plane RMSE (16.24\,mm $\to$ 14.24\,mm) with 15 of 19 captures improving
individually.

\subsection{Validation Summary (Best Configuration)}

\begin{table}[H]
\centering
\caption{Validation statistics for the best configuration
($\lambda = 0.1$, $B_t = 5$\,mm, $B_r = 2.87$\textdegree,
$B_e = 10$\textdegree).}
\label{tab:validation}
\begin{tabular}{@{}lccc@{}}
\toprule
Metric & Uncalibrated & Calibrated & Improvement \\
\midrule
Global RMSE (mm)       & 16.24 & 14.24 & +12.3\% \\
Per-capture RMSE mean  & 14.44 & 12.52 & +13.3\% \\
Per-capture RMSE std   &  7.50 &  6.83 & +8.9\%  \\
Per-capture RMSE max   & 37.82 & 33.65 & +11.0\% \\
Max point error (mm)   & 44.54 & 44.49 & +0.1\%  \\
\bottomrule
\end{tabular}
\end{table}

\subsection{Parameters at Bounds}

Seven of the 14 parameters consistently hit their bounds, indicating
the solver desires corrections larger than what we permit:
\begin{itemize}[nosep]
    \item \code{enc\_tilt\_zero}: $-10.0$\textdegree{} (at bound)
    \item \code{dt\_laser\_x}: $-5.0$\,mm (at bound)
    \item \code{dt\_laser\_y}: $+5.0$\,mm (at bound)
    \item \code{dt\_laser\_z}: $+5.0$\,mm (at bound)
    \item \code{dr\_laser\_roll}, \code{\_pitch}, \code{\_yaw}: all at $\pm 2.87$\textdegree{}
\end{itemize}

The enc\_tilt\_zero absorbs whatever bound is given, which is a hallmark of
a degenerate parameter---it trades off with laser rotation corrections.


% ═══════════════════════════════════════════════════════════════════════
\section{Limitations and Future Work}
\label{sec:future}
% ═══════════════════════════════════════════════════════════════════════

\subsection{Fundamental Limitation: Single-Wall Observability}

The calibration achieves only 7/14 observable parameters because a
single flat wall provides limited geometric diversity. Multiple scan
configurations viewing the same plane can only constrain parameters
that change the \emph{relative} arrangement of points on that plane.
Parameters that rigidly translate or rotate all points together
(such as tilt encoder offset vs.\ laser rotation about the tilt axis)
are perfectly confounded.

\subsection{Recommendations for Improved Calibration}

To break the degeneracy and improve calibration accuracy, the following
strategies should be pursued:

\begin{enumerate}
    \item \textbf{Multiple walls at different orientations.}
          Capture data facing 2--3 walls at different angles. The
          different surface normals provide linearly independent
          constraints that disambiguate currently degenerate parameters.
          A room corner (two perpendicular walls) is the simplest
          improvement.

    \item \textbf{Planar targets at different distances.}
          Walls at varying distances improve sensitivity to translation
          parameters (which scale linearly with range).

    \item \textbf{Geometric features.}
          A corner, edge, or V-shaped target provides line-like
          constraints that constrain both the position and orientation
          of the laser frame simultaneously.

    \item \textbf{Wider joint-space coverage.}
          The current $\pm 25$\textdegree{} pan / $\pm 15$\textdegree{} tilt
          range is limited. If the mechanism allows, sweeping over a
          wider range increases sensitivity to axis alignment errors.

    \item \textbf{Higher-quality lidar sensor.}
          The Hokuyo URG-04LX has specification accuracy of $\pm 60$\,mm
          at ranges beyond 1\,m and inherent noise of 2--5\,mm even at close
          range. A higher-precision sensor (e.g., UTM-30LX) would lower
          the noise floor.

    \item \textbf{Repeated scans and averaging.}
          Taking multiple scans at each configuration and averaging
          reduces random noise. This is a simple change requiring only
          extra acquisition time.

    \item \textbf{Model-based reflectance filtering.}
          The laser's signal quality varies with incidence angle and
          surface material. Filtering by signal intensity (if available)
          could further improve data quality.

    \item \textbf{Reduce the model.}
          Given that only 7 parameters are observable, one could
          fix the remaining 7 at zero and re-solve with only the
          identifiable subset. This eliminates regularization bias
          on the observable parameters. The risk is that the wrong
          parameters may be selected; the SVD analysis should guide
          this selection.

    \item \textbf{Iterative calibration.}
          After an initial calibration updates the URDF, re-collect data
          and solve again. The residual errors after the first round may
          reveal previously masked effects, allowing iterative refinement.
\end{enumerate}

\subsection{Summary of Achievable Accuracy}

\begin{table}[H]
\centering
\caption{Expected accuracy vs.\ calibration environment complexity.}
\label{tab:accuracy_forecast}
\begin{tabular}{@{}lcc@{}}
\toprule
Environment & Expected Obs.\ Rank & Estimated RMSE \\
\midrule
Single wall (current)              & 7/14  & 14--16\,mm \\
Two walls (corner)                 & 10--12/14 & 5--8\,mm \\
Three walls + varying distances    & 13--14/14 & 2--5\,mm \\
High-precision lidar + multi-wall  & 14/14 & $<$2\,mm \\
\bottomrule
\end{tabular}
\end{table}


% ═══════════════════════════════════════════════════════════════════════
\section{Conclusion}
\label{sec:conclusion}
% ═══════════════════════════════════════════════════════════════════════

We presented a complete kinematic calibration pipeline for a pan-tilt
mounted Hokuyo URG-04LX planar LiDAR system. The pipeline includes
automated grid-scan data acquisition, RANSAC-based line fitting,
per-capture line quality filtering, 14-parameter kinematic correction
estimation via Levenberg--Marquardt optimization, SVD-based observability
analysis, and comprehensive validation.

Three key findings emerged:
\begin{enumerate}[nosep]
    \item A single-character Euler convention bug (\code{'xyz'} $\to$
          \code{'XYZ'}) was responsible for 50\% of the initial error,
          highlighting the critical importance of matching rotation
          conventions between the URDF parser and the solver.
    \item The per-capture line quality filter---checking whether individual
          scans conform to a straight line before including them in the
          plane fit---produced a 40\% noise reduction, the single
          largest improvement in the pipeline.
    \item Single-wall calibration fundamentally limits observability to
          7 of 14 parameters, achieving 12.3\% RMSE improvement.
          Breaking through this ceiling requires geometric diversity
          in the calibration environment.
\end{enumerate}

The complete pipeline is implemented as a ROS package
(\code{ptu\_lidar\_calib}) with modular scripts for acquisition,
solving, and validation, ready for deployment on multi-wall environments.


% ═══════════════════════════════════════════════════════════════════════
\appendix
\section{Complete Forward Kinematics Derivation}
\label{app:fk_derivation}
% ═══════════════════════════════════════════════════════════════════════

We detail the full matrix chain for the uncorrected ($\bm{\delta}=\mathbf{0}$)
case. From the URDF:

\paragraph{Step 1: Base to Pan Link.}
\begin{equation}
    \T_{\text{pan}}^{\text{nom}} =
    \begin{bmatrix}
        \mathbf{I}_3 & (0,\ {-0.0285},\ 0.092)^\top \\
        \mathbf{0}^\top & 1
    \end{bmatrix}
\end{equation}
(Identity rotation since RPY = $(0,0,0)$.)

\paragraph{Step 2: Pan Rotation.}
Joint axis is $(0,0,-1)$, so for commanded angle $\theta_p$:
\begin{equation}
    \Rot_{\text{pan}}(\theta_p) = \Rz(-\theta_p) =
    \begin{bmatrix}
        \cos\theta_p & \sin\theta_p & 0 \\
        -\sin\theta_p & \cos\theta_p & 0 \\
        0 & 0 & 1
    \end{bmatrix}
\end{equation}

\paragraph{Step 3: Pan Link to Tilt Link.}
\begin{equation}
    \T_{\text{tilt}}^{\text{nom}} = \T\bigl((0,\ {-0.009},\ 0.022),\; (0,\, 1.57,\, 0)\bigr)
\end{equation}
The RPY $(0,\, 1.57,\, 0)$ applies a $\approx 90$\textdegree{} pitch rotation:
\begin{equation}
    \Rot_y(1.57) \approx
    \begin{bmatrix}
        \cos 1.57 & 0 & \sin 1.57 \\
        0 & 1 & 0 \\
        -\sin 1.57 & 0 & \cos 1.57
    \end{bmatrix}
    \approx
    \begin{bmatrix}
        0 & 0 & 1 \\
        0 & 1 & 0 \\
        -1 & 0 & 0
    \end{bmatrix}
\end{equation}
This rotates the local $z$-axis to align with the global $x$-axis,
so the tilt joint's rotation about local $z$ becomes a rotation about
the (approximate) global $x$-axis---i.e., a tilt.

\paragraph{Step 4: Tilt Rotation.}
Joint axis is $(0,0,1)$, so:
$\Rot_{\text{tilt}}(\theta_t) = \Rz(\theta_t)$.

\paragraph{Step 5: Tilt Link to Laser.}
The composite of two fixed joints:
\begin{equation}
    \T_{\text{laser}}^{\text{nom}} =
    \underbrace{\T\bigl((-0.045, 0, 0),\; (0,0,0)\bigr)}_{\text{hokuyo\_base\_joint}}
    \cdot
    \underbrace{\T\bigl((-0.061, 0.014, 0.040),\; (1.57, 3.14, -1.57)\bigr)}_{\text{hokuyo\_joint}}
\end{equation}

\paragraph{Full Chain.}
\begin{equation}
    \T_{\text{laser}}^{\text{base}} =
    \T_{\text{pan}}^{\text{nom}} \cdot
    \Rz(-\theta_p) \cdot
    \T_{\text{tilt}}^{\text{nom}} \cdot
    \Rz(\theta_t) \cdot
    \T_{\text{laser}}^{\text{nom}}
\end{equation}


\section{Parameter Indices}
\label{app:param_indices}

\begin{table}[H]
\centering
\caption{Mapping from parameter index to physical meaning.}
\label{tab:param_index}
\begin{tabular}{@{}clll@{}}
\toprule
Index & Label & Meaning & Unit \\
\midrule
0  & \code{dt\_pan\_x}   & Pan origin $x$ translation & m \\
1  & \code{dt\_pan\_y}   & Pan origin $y$ translation & m \\
2  & \code{dt\_pan\_z}   & Pan origin $z$ translation & m \\
3  & \code{dr\_pan\_roll}  & Pan origin roll correction & rad \\
4  & \code{dr\_pan\_pitch} & Pan origin pitch correction & rad \\
5  & \code{dr\_pan\_yaw}   & Pan origin yaw correction & rad \\
6  & \code{enc\_pan\_zero}   & Pan encoder zero offset & rad \\
7  & \code{enc\_tilt\_zero}  & Tilt encoder zero offset & rad \\
8  & \code{dt\_laser\_x}  & Laser mount $x$ translation & m \\
9  & \code{dt\_laser\_y}  & Laser mount $y$ translation & m \\
10 & \code{dt\_laser\_z}  & Laser mount $z$ translation & m \\
11 & \code{dr\_laser\_roll}  & Laser mount roll correction & rad \\
12 & \code{dr\_laser\_pitch} & Laser mount pitch correction & rad \\
13 & \code{dr\_laser\_yaw}   & Laser mount yaw correction & rad \\
\bottomrule
\end{tabular}
\end{table}


\end{document}
